\documentclass[a4paper, serif, 10pt]{moderncv}

%% moderncv
% color and style
\moderncvcolor{black}
\moderncvstyle{banking}

%% page layout/margins
\usepackage[top=2.5cm, bottom=2.5cm, left=1.5cm, right=1.5cm]{geometry}

\nopagenumbers{}

%% Babel config for Italian language
% ref:
% - https://latex3.github.io/babel/guides/locale-italian.html
\usepackage[italian]{babel}

%% fontspec
\usepackage{fontspec}

\IfFontExistsTF{Linux Libertine}{
  \setmainfont[Ligatures={TeX, Common}, Numbers=OldStyle]{Linux Libertine}
}{}
\IfFontExistsTF{Linux Libertine O}{
  \setmainfont[Ligatures={TeX, Common}, Numbers=OldStyle]{Linux Libertine O}
}{}

%% CTeX
% CJK support, used to show CJK characters in the resume
%
% - heading=false: disable ctex heading
% - punct=kaiming: use kaiming punctuations styles for CJK
% - scheme=plain: use plain scheme, do not override \normalsize font size
% - space=auto: space settings for CJK characters
%
% CJK fonts are resolved per locale language, e.g. Noto Serif/Sans CJK JP
% for Japanese resumes, with fallback to the next region if not installed
%
% ref:
% - http://ctan.mirrorcatalogs.com/language/chinese/ctex/ctex.pdf
\usepackage[UTF8, heading=false, punct=kaiming, scheme=plain, space=auto]{ctex}

\IfFontExistsTF{Noto Serif CJK SC}{
  \setCJKmainfont{Noto Serif CJK SC}
}{}
\IfFontExistsTF{Noto Sans CJK SC}{
  \setCJKsansfont{Noto Sans CJK SC}
}{}

%% line spacing
\usepackage{setspace}
\setstretch{1.125}

%% URL styling - use normal text instead of monospace
\urlstyle{same}

%% Patch \httplink and \httpslink to handle full URLs
% The original moderncv macros always prepend http:// or https:// to URLs.
% This causes issues when the URL already has a protocol (e.g., https://example.com)
% resulting in malformed URLs like https://https://example.com.
% This patch checks if the URL already contains "://" and uses it directly if so.
% Note: We use \str_set:Nx (x-expansion) to expand macros like \@homepage first.
\ExplSyntaxOn
\str_new:N \g_yamlresume_url_str

\RenewDocumentCommand{\httpslink}{O{}m}{
  \str_set:Nx \g_yamlresume_url_str {#2}
  \str_if_in:NnTF \g_yamlresume_url_str {://}
    { \url{#2} }
    { \url{https://#2} }
}

\RenewDocumentCommand{\httplink}{O{}m}{
  \str_set:Nx \g_yamlresume_url_str {#2}
  \str_if_in:NnTF \g_yamlresume_url_str {://}
    { \url{#2} }
    { \url{http://#2} }
}
\ExplSyntaxOff

\name{Marco Bianchi}{}
\title{Project Manager Senior | Trasformazione digitale e metodologie agili}
\phone[mobile]{+39 347 123 4567}
\email{marco.bianchi@example.it}
\homepage{https://www.marcobianchi-pm.it}

\address{Via Torino 45, Milano, Lombardia, Italia, 20123}{}{}

\extrainfo{{\small \faLinkedin}\ \href{https://www.linkedin.com/in/marcobianchi-pm}{@marcobianchi-pm} {} {} {} • {} {} {} 
{\small \faTwitter}\ \href{https://twitter.com/marcobianchi\_pm}{@marcobianchi\_pm} {} {} {} • {} {} {} 
{\small \faMedium}\ \href{https://medium.com/@marcobianchi}{@marcobianchi}}

\begin{document}

\maketitle

\section{Informazioni personali}

\cvline{}{\begin{itemize}
\item Project manager con oltre 8 anni di esperienza nella conduzione di progetti digitali e di trasformazione organizzativa in contesti complessi.
\item Specializzato in metodologie ibride (Agile, Scrum, Waterfall) e nella gestione del ciclo di vita del progetto, dall'analisi dei requisiti al rilascio.
\item Capacità comprovata di coordinare team cross-funzionali, gestire stakeholder e garantire il rispetto di tempi, budget e qualità.
\item Orientato ai risultati, con forte attenzione alla comunicazione, alla gestione del rischio e al miglioramento continuo.
\end{itemize}}
\section{Formazione}

\cventry{set 2013–lug 2015}
        {Laurea magistrale, Ingegneria Gestionale, Punteggio: 110/110 e lode}
        {Politecnico di Milano}
        {\url{https://www.polimi.it}}
        {}
        {\begin{itemize}
\item Formazione avanzata nella pianificazione, esecuzione e controllo di progetti complessi.
\item Approfondimento di metodologie quantitative per il supporto alle decisioni e la gestione delle risorse.
\end{itemize}
\textbf{Corsi}: Gestione dei progetti, Economia e organizzazione aziendale, Analisi dei sistemi complessi, Gestione dell'innovazione, Gestione operativa}

\cventry{set 2010–lug 2013}
        {Laurea triennale, Economia Aziendale, Punteggio: 108/110}
        {Università degli Studi di Bologna}
        {\url{https://www.unibo.it}}
        {}
        {\begin{itemize}
\item Solida base economica e aziendale per la comprensione dei processi organizzativi.
\item Sviluppo di competenze analitiche e di rendicontazione.
\end{itemize}
\textbf{Corsi}: Economia aziendale, Ragioneria generale, Statistica, Diritto privato}
\section{Esperienza professionale}

\cventry{mar 2022–Presente}
        {Senior Project Manager}
        {Fintech Solutions Italia}
        {\url{https://www.fintechsolutionsitalia.it}}
        {}
        {\begin{itemize}
\item Responsabile del portafoglio di progetti digitali per il lancio di nuovi prodotti di pagamento.
\item Coordinamento di team cross-funzionali composti da responsabili di prodotto, sviluppatori, designer e compliance.
\item Definizione di piani di rilascio, budget e piani di mitigazione dei rischi, con rendicontazione periodica al comitato esecutivo.
\item Introduzione di pratiche agili e metriche di produttività, con riduzione del 25\% del tempo di immissione sul mercato.
\end{itemize}
\textbf{Parole chiave}: Gestione del portafoglio, Metodologie agili, Trasformazione digitale, Gestione del rischio}

\cventry{gen 2019–feb 2022}
        {Project Manager}
        {Digital Retail Group}
        {\url{https://www.digitalretailgroup.it}}
        {}
        {\begin{itemize}
\item Gestione di progetti e-commerce e CRM per marchi del gruppo, con team distribuiti in Italia e all'estero.
\item Pianificazione delle attività, allocazione delle risorse e monitoraggio dei KPI di progetto.
\item Facilitazione di laboratori con stakeholder interni ed esterni per la definizione dei requisiti.
\item Implementazione di un modello di governance basato su Scrum e Kanban.
\end{itemize}
\textbf{Parole chiave}: E-commerce, CRM, Scrum, Kanban, Governance}

\cventry{set 2015–dic 2018}
        {Junior Project Manager}
        {Consulenza Innovazione S.p.A.}
        {\url{https://www.consulenzainnovazione.it}}
        {}
        {\begin{itemize}
\item Supporto alla direzione progetti per clienti nei settori manifatturiero e servizi.
\item Redazione di piani di progetto, analisi dei requisiti e documentazione di avanzamento.
\item Monitoraggio di tempi, costi e qualità, con aggiornamento del registro dei rischi.
\item Partecipazione a gare e presentazioni commerciali per nuovi incarichi.
\end{itemize}
\textbf{Parole chiave}: Pianificazione, Registro dei rischi, Documentazione, Rendicontazione}
\section{Lingue}

\cvline{Italiano}{Madrelingua o bilingue \hfill \textbf{Parole chiave}: Madrelingua}
\cvline{Inglese}{Competenza professionale completa \hfill \textbf{Parole chiave}: IELTS 7.5, Business English}
\cvline{Spagnolo}{Competenza professionale limitata \hfill \textbf{Parole chiave}: DELE B1}
\section{Competenze}

\cvline{Gestione del progetto}{Esperto \hfill \textbf{Parole chiave}: Agile, Scrum, Kanban, Waterfall, PRINCE2}
\cvline{Pianificazione e controllo}{Avanzato \hfill \textbf{Parole chiave}: Microsoft Project, Jira, Asana, Trello, Smartsheet}
\cvline{Analisi dei dati}{Avanzato \hfill \textbf{Parole chiave}: Excel avanzato, Power BI, SQL, Tableau}
\cvline{Comunicazione e leadership}{Esperto \hfill \textbf{Parole chiave}: Negoziazione, Facilitazione, Mentoraggio, Presentazioni efficaci}
\cvline{Gestione del rischio}{Avanzato \hfill \textbf{Parole chiave}: Registro dei rischi, Mitigazione, Audit, Conformità}
\section{Premi}

\cventry{nov 2023}
        {Associazione Italiana Project Management}
        {Premio Eccellenza PMI 2023}
        {}
        {}
        {Riconoscimento per la gestione esemplare di un programma di trasformazione digitale con impatto misurabile su efficienza e soddisfazione del cliente.}

\cventry{giu 2021}
        {Digital Retail Group}
        {Miglior leader di team agile}
        {}
        {}
        {Premio interno per la capacità di guidare team distribuiti e raggiungere gli obiettivi di progetto con elevati standard di qualità.}
\section{Certificazioni}

\cventry{mag 2021}
        {Project Management Institute}
        {Project Management Professional (PMP)}
        {\url{https://www.pmi.org/certifications/project-management-pmp}}
        {}
        {}

\cventry{ott 2019}
        {Scrum Alliance}
        {Certified ScrumMaster (CSM)}
        {\url{https://www.scrumalliance.org/get-certified/scrum-master-track/certified-scrummaster}}
        {}
        {}

\cventry{feb 2020}
        {PeopleCert}
        {ITIL 4 Foundation}
        {\url{https://www.peoplecert.org/browse-certifications/it-governance-and-service-management/ITIL-1/itil-4-foundation}}
        {}
        {}
\section{Pubblicazioni}

\cventry{set 2022}
        {Gestione agile dei progetti complessi: un modello ibrido}
        {Rivista Italiana di Project Management}
        {\url{https://www.rivistaitalianapm.it/articoli/modello-ibrido}}
        {}
        {\begin{itemize}
\item Articolo dedicato all'integrazione tra framework agili e approcci tradizionali.
\item Analisi di casi studio in ambito finanziario e retail.
\item Proposta di un modello di governance per programmi multi-progetto.
\end{itemize}}
\section{Progetti}

\cventry{apr 2022–mag 2023}
        {Progetto di digitalizzazione del processo di onboarding clienti per ridurre i tempi di attivazione e migliorare l'esperienza utente.}
        {Piattaforma di onboarding digitale}
        {\url{https://www.fintechsolutionsitalia.it/progetti/onboarding-digitale}}
        {}
        {\begin{itemize}
\item Coordinamento di un team di 12 persone tra prodotto, sviluppo, design e compliance.
\item Definizione di requisiti funzionali e flussi utente in collaborazione con il business.
\item Rilascio della piattaforma con riduzione del 40\% dei tempi di onboarding.
\item Gestione del budget di 600.000 euro e rendicontazione mensile agli stakeholder.
\end{itemize}
\textbf{Parole chiave}: Onboarding digitale, UX, Conformità, Gestione del budget}

\cventry{mar 2019–dic 2021}
        {Programma pluriennale di migrazione dei sistemi ERP e integrazione dei processi amministrativi, della catena di fornitura e delle vendite.}
        {Programma di trasformazione ERP}
        {\url{https://www.digitalretailgroup.it/progetti/trasformazione-erp}}
        {}
        {\begin{itemize}
\item Pianificazione e coordinamento di 5 flussi di lavoro con team interni e fornitori esterni.
\item Gestione del cambiamento organizzativo e formazione di oltre 200 utenti.
\item Raggiungimento degli obiettivi di avvio operativo entro tempi e budget previsti.
\item Implementazione di un cruscotto di monitoraggio dei KPI dopo l'avvio.
\end{itemize}
\textbf{Parole chiave}: ERP, Gestione del cambiamento, Formazione, KPI}
\section{Interessi}

\cvline{Sport}{Corsa, Ciclismo, Nuoto}
\cvline{Lettura}{Saggistica, Romanzi storici}
\cvline{Volontariato}{Mentoraggio, Formazione}
\section{Volontariato}

\cventry{gen 2021–Presente}
        {Mentore volontario}
        {Associazione Italiana Project Management}
        {\url{https://www.aipm.it}}
        {}
        {\begin{itemize}
\item Attività di mentore per giovani professionisti e studenti che si avvicinano alla gestione dei progetti.
\item Organizzazione di laboratori su metodologie agili e pianificazione.
\item Contributo alla comunità con articoli e incontri tematici.
\end{itemize}}
    

\cventry{set 2018–dic 2020}
        {Formatore volontario}
        {Fondazione Formazione e Lavoro}
        {\url{https://www.fondazioneformazionelavoro.it}}
        {}
        {\begin{itemize}
\item Docenza in corsi base di gestione del progetto per persone in cerca di occupazione.
\item Preparazione di materiali didattici e simulazioni pratiche.
\item Tutoraggio individuale per l'inserimento lavorativo.
\end{itemize}}
    
\end{document}